% GENERATED FILE. Edit canonical lessons, then run npm run generate:books.
% canonical-source-hash: fnv1a64:24c31266101b3b3e
% canonical-lessons: GE-C441-fassade, GE-C441-schaufenster, GE-C441-dachrinne, GE-C441-markthalle, GE-C441-saeule

\chapter{Häuser und Läden}
\label{ch:ge-hauser-und-laden}
\hlchaptermodality{441}

\begin{chapteropening}
\textbf{By the end of this chapter:} \emph{I can describe the fronts of buildings: the façade, the shop window, the gutter, the covered market and the pillar.}

Complete the last lesson of chapter 441: I can describe the fronts of buildings: the façade, the shop window, the gutter, the covered market and the pillar.
\end{chapteropening}

\section[die Fassade]{die Fassade — the façade, the front (of a building)}

\label{lesson:GE-C441-fassade}



\begin{quote}
Before the new one: say the German for the cobbles, then the German for the fountain.
\end{quote}

\subsection*{You'll want to know: die Fassade}
\textbf{die Fassade} — \textquotedblleft{}the façade, the front (of a building)\textquotedblright{}.

\begin{quote}
Die Fassade vom Rathaus ist frisch gestrichen.
\end{quote}

\begin{quote}
\emph{The front of the town hall has been freshly painted.}
\end{quote}

Plural: \emph{die Fassaden}. Figurative: \emph{Das ist nur Fassade} — it is all for show.

\begin{cousinweb}
\emph{Fassade} was borrowed from French \emph{façade} (EtymDB), which came from Italian \emph{facciata}, built on \emph{faccia} 'face', from Latin \emph{faciēs} (Lewis \& Short).
\end{cousinweb}

\begin{grammarlens}[title={What you've built}]
One more word for telling, explaining and describing at B1.
\end{grammarlens}

\subsection*{Your turn}
\begin{itemize}
\raggedright
  \item \emph{Say it:} \emph{die Fassade}
  \item \emph{Say it:} \emph{die Fassade}, once more
  \item \emph{Say it:} a sentence of your own with \emph{die Fassade}
  \item \emph{Recall:} say the German for the cobbles, then the German for the fountain, then say \emph{die Fassade} again
\end{itemize}

\begin{tcolorbox}[breakable,colback=teal!4,colframe=teal!35!black,title={Before you move on}]
What does die Fassade mean? (\textquotedblleft{}The façade, the front (of a building)\textquotedblright{}.) Say it once more.
\end{tcolorbox}
\section[das Schaufenster]{das Schaufenster — the shop window}

\label{lesson:GE-C441-schaufenster}



\begin{quote}
Before the new one: say the German for the fountain, then the German for the façade.
\end{quote}

\subsection*{You'll want to know: das Schaufenster}
\textbf{das Schaufenster} — \textquotedblleft{}the shop window\textquotedblright{}.

\begin{quote}
Im Schaufenster stehen schöne Schuhe.
\end{quote}

\begin{quote}
\emph{There are nice shoes in the shop window.}
\end{quote}

Same form in the plural. \emph{Einen Schaufensterbummel machen} — to go window-shopping.

\begin{cousinweb}
\emph{Schaufenster} is built from \emph{schauen} 'to look', inherited from Old High German \emph{scouwōn} (EtymDB), and \emph{Fenster} 'window', an early borrowing from Latin \emph{fenestra} (Lewis \& Short).
\end{cousinweb}

\begin{grammarlens}[title={What you've built}]
One more word for telling, explaining and describing at B1.
\end{grammarlens}

\subsection*{Your turn}
\begin{itemize}
\raggedright
  \item \emph{Say it:} \emph{das Schaufenster}
  \item \emph{Say it:} \emph{das Schaufenster}, once more
  \item \emph{Say it:} a sentence of your own with \emph{das Schaufenster}
  \item \emph{Recall:} say the German for the fountain, then the German for the façade, then say \emph{das Schaufenster} again
\end{itemize}

\begin{tcolorbox}[breakable,colback=teal!4,colframe=teal!35!black,title={Before you move on}]
What does das Schaufenster mean? (\textquotedblleft{}The shop window\textquotedblright{}.) Say it once more.
\end{tcolorbox}
\section[die Dachrinne]{die Dachrinne — the gutter (on a roof)}

\label{lesson:GE-C441-dachrinne}



\begin{quote}
Before the new one: say the German for the façade, then the German for the shop window.
\end{quote}

\subsection*{You'll want to know: die Dachrinne}
\textbf{die Dachrinne} — \textquotedblleft{}the gutter (on a roof)\textquotedblright{}.

\begin{quote}
Die Dachrinne ist voll mit Laub.
\end{quote}

\begin{quote}
\emph{The gutter is full of leaves.}
\end{quote}

Plural: \emph{die Dachrinnen}.

\begin{cousinweb}
\emph{Dachrinne} is built from \emph{Dach} 'roof' and \emph{Rinne} 'channel', both inherited: \emph{Rinne} goes back to Old High German \emph{rinna}, from the verb for running water (EtymDB).
\end{cousinweb}

\begin{grammarlens}[title={What you've built}]
One more word for telling, explaining and describing at B1.
\end{grammarlens}

\subsection*{Your turn}
\begin{itemize}
\raggedright
  \item \emph{Say it:} \emph{die Dachrinne}
  \item \emph{Say it:} \emph{die Dachrinne}, once more
  \item \emph{Say it:} a sentence of your own with \emph{die Dachrinne}
  \item \emph{Recall:} say the German for the façade, then the German for the shop window, then say \emph{die Dachrinne} again
\end{itemize}

\begin{tcolorbox}[breakable,colback=teal!4,colframe=teal!35!black,title={Before you move on}]
What does die Dachrinne mean? (\textquotedblleft{}The gutter (on a roof)\textquotedblright{}.) Say it once more.
\end{tcolorbox}
\section[die Markthalle]{die Markthalle — the covered market, the market hall}

\label{lesson:GE-C441-markthalle}



\begin{quote}
Before the new one: say the German for the shop window, then the German for the gutter.
\end{quote}

\subsection*{You'll want to know: die Markthalle}
\textbf{die Markthalle} — \textquotedblleft{}the covered market, the market hall\textquotedblright{}.

\begin{quote}
In der Markthalle kauft sie Käse und Obst.
\end{quote}

\begin{quote}
\emph{She buys cheese and fruit in the covered market.}
\end{quote}

Plural: \emph{die Markthallen}.

\begin{cousinweb}
\emph{Markthalle} is built from \emph{Markt} 'market', which goes back through Old High German \emph{marchāt} to Latin \emph{mercātus} (EtymDB), and \emph{Halle} 'hall', an inherited Germanic word.
\end{cousinweb}

\begin{grammarlens}[title={What you've built}]
One more word for telling, explaining and describing at B1.
\end{grammarlens}

\subsection*{Your turn}
\begin{itemize}
\raggedright
  \item \emph{Say it:} \emph{die Markthalle}
  \item \emph{Say it:} \emph{die Markthalle}, once more
  \item \emph{Say it:} a sentence of your own with \emph{die Markthalle}
  \item \emph{Recall:} say the German for the shop window, then the German for the gutter, then say \emph{die Markthalle} again
\end{itemize}

\begin{tcolorbox}[breakable,colback=teal!4,colframe=teal!35!black,title={Before you move on}]
What does die Markthalle mean? (\textquotedblleft{}The covered market, the market hall\textquotedblright{}.) Say it once more.
\end{tcolorbox}
\section[die Säule]{die Säule — the pillar}

\label{lesson:GE-C441-saeule}



\begin{quote}
Before the new one: say the German for the gutter, then the German for the covered market.
\end{quote}

\subsection*{You'll want to know: die Säule}
\textbf{die Säule} — \textquotedblleft{}the pillar\textquotedblright{}.

\begin{quote}
Vor dem Theater stehen sechs hohe Säulen.
\end{quote}

\begin{quote}
\emph{There are six tall pillars in front of the theatre.}
\end{quote}

Plural: \emph{die Säulen}.

\begin{cousinweb}
\emph{Säule} is inherited from Old High German \emph{sūl} 'pillar'; EtymDB lists Gothic \emph{sauls} 'pillar' among its cousins.
\end{cousinweb}

\begin{grammarlens}[title={What you've built}]
One more word for telling, explaining and describing at B1.
\end{grammarlens}

\subsection*{Your turn}
\begin{itemize}
\raggedright
  \item \emph{Say it:} \emph{die Säule}
  \item \emph{Say it:} \emph{die Säule}, once more
  \item \emph{Say it:} a sentence of your own with \emph{die Säule}
  \item \emph{Recall:} say the German for the gutter, then the German for the covered market, then say \emph{die Säule} again
\end{itemize}

\begin{tcolorbox}[breakable,colback=teal!4,colframe=teal!35!black,title={Before you move on}]
What does die Säule mean? (\textquotedblleft{}The pillar\textquotedblright{}.) Say it once more.
\end{tcolorbox}
